%% cvforge-template v1
%% sections: experience, projects, skills, education
%% supports: person.phone, person.links, experience.location, projects.url
%% max-bullets-per-role: 5
%
% AltaCV (LPPL 1.3 or later, LianTze Lim), converted to the cvforge contract. Class, packages,
% fonts and two-column layout are the original sample's. Sample-only sections with no profile
% equivalent (photo, philosophy quote, achievements, languages, referees, wheel chart,
% publications) are removed; "Strengths" holds the skills as the sample's \cvtag list did.
%%%%%%%%%%%%%%%%%
% This is an sample CV template created using altacv.cls
% (v1.7.4, 30 July 2025) written by LianTze Lim (liantze@gmail.com). Compiles with pdfLaTeX, XeLaTeX and LuaLaTeX.
%
%% It may be distributed and/or modified under the
%% conditions of the LaTeX Project Public License, either version 1.3
%% of this license or (at your option) any later version.
%% The latest version of this license is in
%%    http://www.latex-project.org/lppl.txt
%% and version 1.3 or later is part of all distributions of LaTeX
%% version 2003/12/01 or later.
%%%%%%%%%%%%%%%%

%% Use the "normalphoto" option if you want a normal photo instead of cropped to a circle
% \documentclass[10pt,a4paper,withhyper,normalphoto]{altacv}

\documentclass[10pt,a4paper,withhyper]{altacv}
%% AltaCV uses the fontawesome5 and simpleicons packages.
%% See http://texdoc.net/pkg/fontawesome5 and http://texdoc.net/pkg/simpleicons for full list of symbols.

% Change the page layout if you need to
\geometry{left=1.25cm,right=1.25cm,top=1.5cm,bottom=1.5cm,columnsep=1.2cm}

% The paracol package lets you typeset columns of text in parallel
\usepackage{paracol}

% Change the font if you want to, depending on whether
% you're using pdflatex or xelatex/lualatex
% WHEN COMPILING WITH XELATEX PLEASE USE
% xelatex -shell-escape -output-driver="xdvipdfmx -z 0" sample.tex
\iftutex
  % If using xelatex or lualatex:
  \setmainfont{Roboto Slab}
  \setsansfont{Lato}
  \renewcommand{\familydefault}{\sfdefault}
\else
  % If using pdflatex:
  \usepackage[rm]{roboto}
  \usepackage[defaultsans]{lato}
  % \usepackage{sourcesanspro}
  \renewcommand{\familydefault}{\sfdefault}
\fi

% Change the colours if you want to
\definecolor{SlateGrey}{HTML}{2E2E2E}
\definecolor{LightGrey}{HTML}{666666}
\definecolor{DarkPastelRed}{HTML}{450808}
\definecolor{PastelRed}{HTML}{8F0D0D}
\definecolor{GoldenEarth}{HTML}{E7D192}
\colorlet{name}{black}
\colorlet{tagline}{PastelRed}
\colorlet{heading}{DarkPastelRed}
\colorlet{headingrule}{GoldenEarth}
\colorlet{subheading}{PastelRed}
\colorlet{accent}{PastelRed}
\colorlet{emphasis}{SlateGrey}
\colorlet{body}{LightGrey}

% Change some fonts, if necessary
\renewcommand{\namefont}{\Huge\rmfamily\bfseries}
\renewcommand{\personalinfofont}{\footnotesize}
\renewcommand{\cvsectionfont}{\LARGE\rmfamily\bfseries}
\renewcommand{\cvsubsectionfont}{\large\bfseries}


% Change the bullets for itemize and rating marker
% for \cvskill if you want to
\renewcommand{\cvItemMarker}{{\small\textbullet}}
\renewcommand{\cvRatingMarker}{\faCircle}
% ...and the markers for the date/location for \cvevent
% \renewcommand{\cvDateMarker}{\faCalendar*[regular]}
% \renewcommand{\cvLocationMarker}{\faMapMarker*}


% If your CV/résumé is in a language other than English,
% then you probably want to change these so that when you
% copy-paste from the PDF or run pdftotext, the location
% and date marker icons for \cvevent will paste as correct
% translations. For example Spanish:
% \renewcommand{\locationname}{Ubicación}
% \renewcommand{\datename}{Fecha}


%% Use (and optionally edit if necessary) this .tex if you
%% want to use an author-year reference style like APA(6)
%% for your publication list
% \input{pubs-authoryear.cfg}

%% Use (and optionally edit if necessary) this .tex if you
%% want an originally numerical reference style like IEEE
%% for your publication list
% \input{pubs-num.cfg}  % cvforge: publications are not in the profile

%% sample.bib contains your publications
% \addbibresource{sample.bib}  % cvforge: publications are not in the profile

\begin{document}
\name{Tomasz Lindqvist}
\tagline{Backend Developer}
%% cvforge: the sample photo (\photoR) is removed; there is no photo in the profile.

\personalinfo{%
  \email{tomasz.lindqvist@example.com}
  \phone{+48 500 000 000}
  \location{Gdansk, PL}
  \homepage{github.com/tlindqvist-example}
  \homepage{linkedin.com/in/tlindqvist-example}
}

\makecvheader

\columnratio{0.6}

\begin{paracol}{2}
\cvsection{Experience}

\cvevent{Senior Backend Engineer}{Brightharbor Logistics}{Mar 2023 -- Present}{Gdansk, PL}
\begin{itemize}
\item Designed a public REST API in Go with OpenAPI 3 specs, serving 40M requests a month to 300 partner integrations. % id:exp-01.b1
\item Introduced gRPC between internal services, reducing p99 latency from 480 ms to 90 ms. % id:exp-01.b3
\item Migrated 22 services from VMs to Kubernetes on Terraform-managed AWS, cutting hosting cost 35\%. % id:exp-01.b2
\item Built Prometheus alerting and SLO dashboards that cut mean time to recovery from 3 hours to 25 minutes. % id:exp-01.b4
\end{itemize}

\divider

\cvevent{Backend Developer}{Copperline Payments}{Jun 2020 -- Feb 2023}{Warsaw, PL}
\begin{itemize}
\item Built an idempotent payments API in Python and PostgreSQL processing 1.2M transactions a day. % id:exp-02.b1
\item Replaced nightly batch jobs with Kafka event streams, delivering settlement data 6 hours sooner. % id:exp-02.b2
\item Added Redis caching to the account lookup path, cutting database load 60\%. % id:exp-02.b3
\item Containerised 9 services with Docker and wrote the local setup used by 25 engineers. % id:exp-02.b4
\end{itemize}

\divider

\cvevent{Junior Software Developer}{Fernwood Analytics}{Jul 2018 -- May 2020}{Warsaw, PL}
\begin{itemize}
\item Wrote 14 REST endpoints in Python for a reporting product used by 120 clients. % id:exp-03.b1
\item Tuned slow PostgreSQL queries, bringing the 10 heaviest reports under 2 seconds. % id:exp-03.b2
\item Automated test database provisioning with Docker, saving each developer 30 minutes a day. % id:exp-03.b3
\end{itemize}

\cvsection{Projects}

\cvevent{queuectl}{github.com/tlindqvist-example/queuectl}{}{}
\begin{itemize}
\item Open-source Go CLI for inspecting Kafka consumer lag, with 1,400 GitHub stars. % id:prj-01.b1
\end{itemize}

\switchcolumn

\cvsection{Strengths}

\cvtag{Go}
\cvtag{Python}
\cvtag{REST APIs}
\cvtag{PostgreSQL}
\cvtag{Apache Kafka}
\cvtag{OpenAPI}
\cvtag{Kubernetes}
\cvtag{gRPC}
\cvtag{Terraform}
\cvtag{AWS}
\cvtag{Redis}
\cvtag{Docker}
\cvtag{Prometheus}

\medskip

\cvsection{Education}

\cvevent{B.Sc. Computer Science}{Lakeside Technical University}{Oct 2014 -- Jun 2018}{}

\end{paracol}

\end{document}
